\begin{abstract}
We audit seven public phishing URL corpora, whose two classes come from different collection pipelines, for presentation artefacts in recorded strings, sampling artefacts such as homepage-only legitimate sets, and registered-domain leakage. On PhiUSIIL a classifier restricted to five formatting flags and the scheme type reaches an MCC of \WPhiUSIILRandomPresentationHttps{} (\WPhiUSIILRandomPresentation{} without it), and this contrast persists among homepages alone; on a PhishTank and Alexa construction the scheme flag alone separates the classes. Canonicalising the strings lowers a character n-gram detector by \DRawVsCanonPhiUSIIL{} MCC on PhiUSIIL and \DRawVsCanonKaggleMU{} on the Kaggle malicious URL collection, and by less than 0.01 on average on the four other published corpora; a character CNN shows the same two large drops. Under registered-domain grouping a host-only model on ISCX-URL2016 falls from \WISCXURLRandomCharHost{} to \WISCXURLDomainCharHost{} because its legitimate URLs come from \RegLegitISCXURL{} domains. The Kaggle collection contains every legitimate ISCX-URL2016 URL and labels \XConfAPKaggleMUPhishStorm{} PhishStorm legitimate URLs as phishing. Across corpora, after purging overlapping test URLs, canonicalisation raises mean ROC-AUC from \XUrlCharRawMeanAUC{} to \XUrlCharCanonMeanAUC, mainly where corpora use opposite conventions, while formatting flags alone transfer at chance on average (pair AUC \XUrlPresentationMinAUC{} to \XUrlPresentationMaxAUC). Thresholds set on validation data for a 0.1\% false-positive rate reach a median of \XVFprCharCanon\% on other corpora, leaving a median precision of \XVPrecOneCharCanon{} at 1\% prevalence. We propose an eight-item reporting checklist.
\end{abstract}

\begin{keyword}
phishing detection \sep URL classification \sep dataset bias \sep shortcut learning \sep data leakage \sep benchmark audit \sep network security
\end{keyword}
